\chapter{Introducción}

\section{Contexto}

Los paradigmas de toma de decisiones humanas se estudian habitualmente mediante
modelos matemáticos, experimentos conductuales y simulaciones. Sin embargo,
pasar de una descripción científica a un modelo ejecutable sigue siendo una
tarea manual: exige buscar literatura, identificar variables, escoger una
formulación, adaptarla a un entorno concreto y escribir código correcto.

Este trabajo aborda esa brecha desde la ingeniería de agentes inteligentes. La
fase desarrollada por Pablo Pazos Parada construye un pipeline capaz de partir
de una descripción en lenguaje natural y producir modelos de decisión
ejecutables. La salida no es un texto explicativo, sino código Python probado y
preparado para integrarse en el laboratorio virtual de simulación.

\begin{figure}[H]
\centering
\begin{verbatim}
Descripcion del problema
        |
        v
Investigacion cientifica
        |
        v
Formalizacion matematica
        |
        v
Adaptacion al entorno
        |
        v
Modelo Python probado
\end{verbatim}
\caption{Transformación principal de la fase 1.}
\end{figure}

\section{Motivación}

El proyecto completo, compartido con la fase de Juan Freire Alvarez, persigue un
laboratorio virtual para simular, observar y analizar agentes basados en modelos
de decisión. Dentro de ese objetivo, esta memoria se centra en la parte que
genera dichos modelos.

La motivación técnica es doble. En primer lugar, los modelos de lenguaje pueden
ayudar a interpretar literatura científica y sintetizar código, pero necesitan
control, validación y trazabilidad. En segundo lugar, un sistema de investigación
no debe empezar desde cero en cada ejecución: los paradigmas, formulaciones,
decisiones de diseño y modelos aceptados deben convertirse en memoria reutilizable.

\section{Objetivos}

El objetivo general es diseñar e implementar un sistema de agentes que genere
modelos de decisión ejecutables a partir de problemas descritos en lenguaje
natural.

Los objetivos específicos son:

\begin{itemize}
\item Separar el proceso en etapas especializadas: investigación, formalización,
razonamiento de implementación y construcción.
\item Mantener revisión humana entre etapas para evitar que errores tempranos se
propaguen al código final.
\item Persistir artefactos, trazas y estado de ejecución para poder auditar y
reanudar el pipeline.
\item Crear una memoria de conocimiento que combine base de datos relacional,
almacenamiento de objetos, grafo de conocimiento y recuperación vectorial.
\item Generar modelos Python compatibles con un contrato mínimo de simulación:
\texttt{decide}, \texttt{update} y \texttt{get\_state}.
\item Validar los modelos generados mediante pruebas ejecutables.
\end{itemize}

\section{Alcance}

Esta memoria describe la fase 1 del sistema. Quedan dentro del alcance:

\begin{itemize}
\item El enrutador que coordina el pipeline.
\item Los agentes Researcher, Formalizer, Reasoner, Builder y MemoryAgent.
\item La integración con infraestructura compartida: PostgreSQL, MinIO, Neo4j,
Qdrant, embeddings y reranking.
\item El contrato de salida de los modelos generados.
\item Las decisiones de diseño tomadas para controlar calidad, trazabilidad y
memoria.
\end{itemize}

Quedan fuera del alcance principal, aunque se mencionan cuando son necesarios:
la interfaz web, el motor de simulación completo, los agentes de análisis e
informes de la segunda fase y la evaluación empírica de paradigmas concretos.

\section{Estructura de la memoria}

El capítulo 2 resume el estado del arte relevante. El capítulo 3 presenta la
arquitectura global. El capítulo 4 profundiza en el flujo de agentes. El
capítulo 5 describe la memoria y el grafo de conocimiento. El capítulo 6 recoge
detalles de implementación y validación. Finalmente, el capítulo 7 presenta las
conclusiones y líneas futuras.
